\chapter{Introduction}\label{s1:chapter}

\section{Yang--Mills theory and its lattice formulation}\label{s1:background}

Yang and Mills introduced non-Abelian gauge theories in 1954 \cite{YangMills1954}. In the Euclidean formulation of the quantum theory, one would like to define a probability measure on connections $A$ on a three- or four-dimensional manifold, formally given by
\begin{equation}\label{s1:eq:formal}
 d\mu(A)=Z^{-1}\exp\Bigl(-\frac{1}{2g^2}\int|F_A|^2\,dx\Bigr)\,\mathcal DA.
\end{equation}
Here $F_A$ is the curvature of $A$, $g$ is the coupling constant, and $\mathcal DA$ stands for a Lebesgue measure on the infinite-dimensional space of connections, which does not exist. The measure should be invariant under gauge transformations, so it is really a measure on gauge orbits. The natural observables are gauge-invariant functions, such as traces of holonomies around closed loops, which are called Wilson loops.

Wilson's lattice regularization \cite{Wilson1974} replaces the connection by variables $U_e$ in the gauge group, one for each link $e$ of a lattice with mesh $a$. The action becomes the sum of $1-\tr U_p/2$ over the plaquettes $p$ of the lattice, where $U_p$ is the product of the link variables around $p$. For every $a>0$ this defines a probability measure that is exactly invariant under lattice gauge transformations. Osterwalder and Seiler established reflection positivity and other basic properties of these measures \cite{OsterwalderSeiler1978}; see also Seiler's monograph \cite{Seiler1982}. The central problem is to remove the regularization: one must show that suitable observables have limits as $a\to0$, with the bare coupling tuned appropriately, and describe the limit. In four dimensions this is part of the problem described by Jaffe and Witten \cite{JaffeWitten2006}; see \cite{Douglas2026} for a recent review by Douglas.

In this paper we consider the gauge group $\SU$ in three dimensions, on a periodic torus $\Torus=(\R/L\Z)^3$ of fixed side $L$. In three dimensions the coupling $g_3^2$ in \eqref{s1:eq:formal} has the dimension of inverse length. The dimensionless strength of the interaction at length scale $\varepsilon$ is $g_3^2\varepsilon$. In the sense of power counting, the theory becomes weakly coupled at short distances; it is superrenormalizable. The bare lattice coupling at mesh $a$ is $g^2=g_3^2a$. With $N=L/a$ lattice sites along each side, this means that
\begin{equation}\label{s1:eq:coupling}
 g_N^2=\frac{\gamma}{N},\qquad \gamma=g_3^2L,
\end{equation}
where the dimensionless number $\gamma$ is the coupling measured in units of the size of the torus.

Three-dimensional Yang--Mills theory is of independent physical interest. It describes the magnetic sector of four-dimensional gauge theory at high temperature and long distances; see Ginsparg \cite{Ginsparg1980} and Appelquist and Pisarski \cite{AppelquistPisarski1981}, and also Linde \cite{Linde1980} and Jackiw and Templeton \cite{JackiwTempleton1981} on the infrared problems of its perturbative treatment. It is also expected to have a mass gap and to confine. Feynman gave qualitative arguments for both \cite{Feynman1981}. Karabali, Kim, and Nair developed a gauge-invariant Hamiltonian analysis of the gap and an approximate prediction of the string tension \cite{Karabali1998a,Karabali1998b}, and Teper computed the spectrum and string tensions numerically, with continuum extrapolations \cite{Teper1999}. Our results concern only the ultraviolet problem in a fixed finite volume.

Even in three dimensions, the ultraviolet problem is not only a question of bounds. In the Gaussian (Maxwell) approximation to \eqref{s1:eq:formal}, after gauge fixing and with the constant modes treated separately, a connection in three dimensions is a random distribution of H\"older--Besov regularity $\alpha$ for every $\alpha<-1/2$, and no better. Holonomies along general fixed loops then require an additional regularization \cite{CCState2023v4,CCHeat2022v4}. A convenient way to obtain gauge-invariant observables with limits is to smooth the field first. The gradient flow of the Yang--Mills action, $\partial_tA=-D_A^*F_A$, smooths connections at length scale $\sqrt t$ and commutes with gauge transformations. Its lattice version is the gradient flow of the Wilson action. Narayanan and Neuberger studied continuum limits of smeared Wilson loops numerically \cite{NarayananNeuberger2006}, and L\"uscher developed this lattice flow systematically under the name Wilson flow \cite{Luscher2010a,Luscher2010WilsonFlow}. In \cite{LuscherWeisz2011}, L\"uscher and Weisz showed that in four dimensions the correlation functions of the flowed field are finite to all orders in renormalized perturbation theory. Wilson loops of the flowed field at positive flow time are bounded gauge-invariant functions of the lattice configuration, and we use them as our observables.

The geometric study of Yang--Mills connections and their flow has a long history, including work of Atiyah and Bott \cite{AtiyahBott1983} and Donaldson \cite{Donaldson1983,Donaldson1985}. In three dimensions, global existence and regularity for smooth initial data are due to R\aa de \cite{Rade1992}. Charalambous and Gross developed the flow for rough initial data and proposed it as a way to define gauge fields \cite{CharalambousGross2013,CharalambousGross2015,CharalambousGross2017,Gross2022}. In \cite{CCState2023v4}, Cao and Chatterjee made this proposal precise. They defined a state space for three-dimensional Yang--Mills theory, consisting of gauge orbits of trajectories of the Yang--Mills heat flow at positive times. They showed that it carries a natural topology in which Wilson loops of the flowed field are continuous, and they gave a tightness criterion for measures on it. Their arguments rest on Uhlenbeck's compactness theorem \cite{Uhlenbeck1982Connections} and on Sengupta's result that holonomies separate gauge orbits \cite{Sengupta1994}. In \cite{CCHeat2022v4} they also proved local existence of the flow in the DeTurck gauge for a class of random distributional initial data that includes the Gaussian free field. Their decomposition into a rough linear part and a smoother nonlinear part follows the random-data method of Bourgain \cite{Bourgain1996} and of Da Prato and Debussche \cite{DaPratoDebussche2003}. The construction of a Yang--Mills measure on this state space from the lattice theory was left open. Our main result is such a construction.

\section{Previous rigorous results}\label{s1:previous}

We summarize the rigorous results most closely related to ours. The literature is large, and Chatterjee's survey \cite{Chatterjee2019Survey} gives an overview from the point of view of probability.

\paragraph{Abelian theories, Higgs models, and two dimensions.}
In \cite{Gross1983Abelian}, Gross proved that three-dimensional pure $U(1)$ lattice gauge theory converges to the free electromagnetic field, for the Villain action and, in the electric sector, for the Wilson action. Brydges, Fr\"ohlich, and Seiler constructed the Abelian Higgs model in two dimensions \cite{Brydges1979a,Brydges1979b,Brydges1980,Brydges1981}. King constructed the Abelian Higgs model in two and three dimensions by proving convergence of effective actions \cite{King1986a,King1986b}. Chatterjee constructed a Gaussian scaling limit of $\SU$ lattice Yang--Mills theory coupled to a fixed-length Higgs field in the fundamental representation, in every dimension $d\ge2$ \cite{Chatterjee2026HiggsScaling}. The result uses unitary gauge, sends the Higgs length to infinity and the gauge coupling to zero rapidly, and yields independent massive Proca fields; it concerns a different scaling regime from the pure gauge theory considered here. In two dimensions, Yang--Mills theory is exactly solvable. Gross, King, and Sengupta \cite{Gross1989} and Driver \cite{Driver1989} gave probabilistic constructions in the plane, and Driver proved that the lattice Wilson loop expectations converge. Sengupta \cite{Sengupta1992,Sengupta1993,Sengupta1997} and L\'evy \cite{Levy2003,Levy2010} treated compact surfaces, and Fine gave a different geometric approach \cite{Fine1990,Fine1991}. In \cite{Witten1991}, Witten related the weak-coupling limit on surfaces to volumes of moduli spaces of flat connections. In \cite{Chevyrev2019}, Chevyrev realized the two-dimensional measure on the torus as a random distributional connection.

\paragraph{The renormalization group for lattice gauge theories.}
Bałaban developed a renormalization group method for lattice gauge theories in a long series of papers. He proved estimates for lattice Green's functions and propagators \cite{Balaban1983Green,Balaban1984I,Balaban1984II,Balaban1985Background}, introduced a gauge-covariant averaging operation for group-valued fields \cite{Balaban1985Averaging}, and analyzed regular configurations and the variational problem that defines the background field \cite{Balaban1985Regular,Balaban1985Variational}. The culmination is ultraviolet stability on a fixed torus: bounds on renormalized effective densities and suitably normalized partition functions that are uniform in the lattice spacing. He proved it in three dimensions in \cite{Balaban1985UV}, and the later papers on effective actions and large fields developed the four-dimensional program \cite{Balaban1987RGI,Balaban1988RGII,Balaban1989LargeI,Balaban1989LargeII}. His treatment of the complete effective densities also gave convergent renormalization expansions for superrenormalizable gauge models \cite{Balaban1988Convergent}. Bałaban, Imbrie, and Jaffe explained the exact renormalization group for gauge theories and treated Higgs models \cite{Balaban1984d,Balaban1985f,BalabanJaffe1986,Balaban1988c}. Federbush developed a different phase-cell approach to the stability problem \cite{Federbush1986,Federbush1987a,Federbush1987b,FederbushWilliamson1987,Federbush1988,Federbush1990}. Ultraviolet stability is a statement about bounds, and it does not by itself give convergence of observables or a limiting measure.

In \cite{MagnenRivasseauSeneor1993}, Magnen, Rivasseau, and S\'en\'eor developed a constructive analysis of the Schwinger functions of four-dimensional pure $\SU$ Yang--Mills theory with a fixed infrared cutoff, in the trivial topological sector. They work in a regularized axial gauge, with counterterms that restore gauge invariance, and they prove large-field estimates and nonperturbative Slavnov identities. Dimock gave a detailed exposition of Bałaban's method for the scalar $\phi^4_3$ model \cite{Dimock2013v2,Dimock2013Large,Dimock2014Convergence}, and applied related ideas to gauge fixing and to quantum electrodynamics in three dimensions \cite{Dimock2015Axial,Dimock2015b,Dimock2018,Dimock2020}. More broadly, our methods descend from the block-spin renormalization group of Gaw\c{e}dzki and Kupiainen \cite{GawedzkiKupiainen1980,GawedzkiKupiainen1983a,GawedzkiKupiainen1986}. They also use the cluster expansion methods described by Glimm and Jaffe \cite{GlimmJaffe1987}, Brydges \cite{Brydges1986,Brydges2009}, and Rivasseau \cite{Rivasseau1991}. More specifically, we use the convergence criterion of Koteck\'y and Preiss \cite{KoteckyPreiss1986} and the tree estimates going back to Penrose, as explained by Fern\'andez and Procacci \cite{FernandezProcacci2007} and recorded in the graph estimate of Dimock \cite{Dimock2013v2}.

\paragraph{Stochastic quantization.}
A different approach constructs the Yang--Mills measure as the invariant measure of a Langevin dynamics. This idea goes back to Parisi and Wu \cite{ParisiWu1981}, with the gauge-fixing term introduced by Zwanziger \cite{Zwanziger1981}. Using regularity structures \cite{Hairer2014}, Chandra, Chevyrev, Hairer, and Shen constructed a gauge-covariant Langevin dynamics and a state space of gauge orbits in two dimensions \cite{ChandraChevyrevHairerShen2022}. They also constructed local-in-time dynamics for the Yang--Mills--Higgs model in three dimensions \cite{ChandraChevyrevHairerShen2024}. Chevyrev and Shen proved that the Yang--Mills measure on the trivial bundle over the two-dimensional torus is invariant for this dynamics \cite{ChevyrevShen2026Invariant}. They also proved that the mass renormalization of the Yang--Mills field that yields gauge-covariant solutions of the stochastic three-dimensional Yang--Mills--Higgs equations is unique \cite{ChevyrevShen2026Uniqueness}. Bringmann and Cao gave a paracontrolled approach in two dimensions \cite{BringmannCao2026}. For the non-Abelian theory in three dimensions, these results give local dynamics, and they do not yet construct an invariant measure. For the Abelian theory, Chevyrev, Qu, and Shen proved convergence of gauge-fixed three-dimensional lattice Langevin dynamics to the one-form stochastic heat equation for a class of plaquette actions including Wilson, Manton, and Villain actions, under assumptions on the initial data \cite{ChevyrevQuShen2026}. We refer to Chevyrev's survey \cite{Chevyrev2022Survey} for more on this approach.

\paragraph{Lattice estimates.}
Many recent results concern lattice gauge theories at a fixed lattice spacing, for instance at strong coupling or large rank \cite{Chatterjee2019a,ChatterjeeJafarov2016,BasuGanguly2018,Shen2023,CaoNissim2025a,AdhikariCao2025}. Closer to the present work are results on the free energy and on bounds that are uniform in the lattice spacing. In \cite{Chatterjee2016}, Chatterjee identified the leading term of the $U(N)$ lattice Yang--Mills free energy, and in \cite{Brennecke2026}, Brennecke computed its Maxwell constant. Faria da Veiga and O'Carroll proved stability bounds and bounds for generating functions of scaled plaquette fields \cite{FariaVeigaOCarroll2019,FariaVeigaOCarroll2025}. These results control the normalization and local plaquette observables. They do not identify the continuum laws of macroscopic flowed Wilson loops, and asymptotics of the free energy that are extensive in the volume do not by themselves control normalized observables.

\section{Main results}\label{s1:results}

We now describe our results informally; the precise statements are in \cref{s2:chapter}. Let $\mu_{N,\gamma}$ denote the Wilson measure for $\SU$ on the periodic lattice in $\Torus$ with $N^3$ sites and mesh $a_N=L/N$, with the coupling \eqref{s1:eq:coupling}. For a configuration $U$ and a physical time $t>0$, we write $U(t)$ for the Wilson flow of $U$ run for the lattice time $4t/a_N^2$; with this normalization, $U(t)$ approximates the continuum Yang--Mills heat flow at time $t$. Let $\ell$ be a piecewise $C^1$ closed loop in $\Torus$, and let $\ell_N$ be a lattice approximation of $\ell$ of bounded length, in the sense of \cref{s2:def:approx}. The flowed Wilson loop is $\tr\Hol_{\ell_N}U(t)/2$, the normalized trace of the product of the flowed link variables along $\ell_N$.

We fix an odd integer $M\ge3$ and an integer $n$ divisible by four, and consider the lattice sides $N_k=nM^k$ for $k\ge0$. The restriction to such sequences comes from the renormalization group, which divides the lattice into blocks of side $M$.

\begin{itemize}
\item \emph{Existence of the continuum limit} (\cref{s2:thm:main}). Fix $L>0$ and $\gamma>0$. Consider finitely many flowed Wilson loops, for arbitrary piecewise $C^1$ loops, characters of $\SU$, and positive flow times. As $k\to\infty$, their joint laws under $\mu_{N_k,\gamma}$ converge, for every choice of admissible lattice approximations, and the limits do not depend on this choice. All these limits are realized by a single probability measure $\lambda$ on the space $\OrbitSpace$ of gauge orbits of Yang--Mills heat-flow trajectories at positive times. That is, the limiting law of each finite family of flowed Wilson loops is the law of the corresponding continuum holonomies under $\lambda$. The energy at each positive flow time is tight. The measure $\lambda$ depends only on $L$, $\gamma$, and the class of $n$ under multiplication by powers of $M$, and it is invariant under translations, rotations, and reflections of the lattice, and under the center symmetries.
\item \emph{A bound on large fields that is uniform in the coupling} (\cref{s2:thm:large-field}). We carry out the renormalization group analysis on a small-field event, an explicit set of configurations that are close to a gauge transform of the identity at every scale of a multiscale decomposition. Let $u$ be the side of the coarsest lattice in the decomposition. We prove that the complement of the event has probability at most $Cu^3\exp(-cu/\gamma)$, plus a term of order $u^3\exp(-c\log^2(u/\gamma))$. The bound is uniform in the number of renormalization group steps and in $\gamma$ below any fixed ceiling.
\item \emph{Short distances} (\cref{s2:thm:short-distance}). Fix a smooth loop $\ell$ in $\R^3$ and $\sigma>0$, and let $W$ be the flowed Wilson loop of the loop $x+\varepsilon\ell$ at flow time $\sigma\varepsilon^2$. As $\varepsilon\to0$, the normalized deficit $(1-W)/(g_3^2\varepsilon)$ converges in law to $|C|^2/8$, where $C$ is an explicit Gaussian vector given by the free field. This holds at every fixed coupling and uniformly in the position of the loop. Jointly for finitely many fixed loop shapes, the limit is a quadratic function of a Gaussian vector, and it is asymptotically independent of every fixed bounded measurable function on $\OrbitSpace$. This makes precise, for flowed Wilson loops, the statement that the theory is weakly coupled at short distances. In three dimensions this is a consequence of superrenormalizability, rather than of the logarithmic running of the coupling found in four dimensions by Gross and Wilczek \cite{GrossWilczek1973} and Politzer \cite{Politzer1973}.
\item \emph{Weak coupling} (\cref{s2:thm:flat,s2:thm:regular,s2:thm:central}). As $\gamma\to0$, every subsequential limit of the measures $\lambda_\gamma$ is supported on stationary flat connections. In particular, the holonomies $(H_1,H_2,H_3)$ around the three coordinate cycles concentrate on commuting triples. On compact regular patches near the identity, the four normal components of the logarithms of the flowed holonomies have asymptotically Gaussian fluctuations of size $\sqrt\gamma$, with an explicit covariance, in a sense made precise in \cref{s2:thm:regular}. At the central triples the flat connections are singular, and the logarithms of the holonomies live on the scale $\gamma^{1/4}$. Ratios of expectations of compactly supported invariant functions of the rescaled logarithms converge to the corresponding ratios of integrals against $e^{-Q(Y)}\,dY$, where $Q(Y)=\sum_{i<j}|y_i\times y_j|^2/2$ is the constant-mode quartic form. The error is $O(\gamma^{1/4})$ for Lipschitz test functions. The measure $e^{-Q(Y)}\,dY$ has infinite total mass, so this is not a statement about a limiting probability measure. In both regimes we obtain explicit statistics that distinguish $\lambda_\gamma$ from specific Gaussian and Abelian reference measures; the leading coefficient in the regular regime is itself the coefficient of the Gaussian approximation around a regular flat connection.
\end{itemize}

\Cref{s2:chapter} also contains further results, including continuity of $\lambda$ in $\gamma$, with convergence that is uniform on compact sets of couplings, and exact scaling in $L$. We also prove convergence of the energy and loop processes on compact intervals of positive times, and convergence in the Sobolev topology of Cao and Chatterjee. Finally, for $0<\gamma\le1$ and along the subsequences of \cref{s2:rem:ladders}, the limit does not change when the Wilson flow is replaced by the gradient flow of $\sum_p[q(U_p)+c\,q(U_p)^2]$, uniformly for $c$ in a fixed bounded interval of nonnegative numbers.

We emphasize that the main theorem holds at every fixed coupling $\gamma>0$, not only at small coupling. This reflects superrenormalizability: at lattice side $u$ the effective coupling is $\gamma/u$, so we can run the renormalization group from the finest lattice down to a coarse lattice whose side $u$ is large compared to $\gamma$, where the field is still weakly coupled. By the large-field bound, the probability discarded by this procedure tends to zero as $u\to\infty$, uniformly in the lattice spacing. The results on weak coupling concern the opposite limit, $\gamma\to0$ at fixed $u$, taken after the continuum limit.

\section{Relation to previous work}\label{s1:relation}

Compared to Bałaban's three-dimensional theorem \cite{Balaban1985UV}, which bounds renormalized effective densities and suitably normalized partition functions uniformly in the cutoff, our main theorem identifies the limits of normalized observables and constructs a limiting measure. Our proof follows Bałaban's general strategy. We use several of his constructions: the averaging operation, the background field obtained by minimizing the action subject to block constraints, and estimates for the associated propagators. However, we organize the analysis differently. Bałaban decomposes the configurations into small-field and large-field regions at every scale. Instead, we condition on a single small-field event. We bound the probability of its complement directly under the Wilson measure, using the reflection positivity of Osterwalder and Seiler \cite{OsterwalderSeiler1978} and the chessboard estimates of Fr\"ohlich, Israel, Lieb, and Simon \cite{FILS1978,FILS1980}, in the form given by Biskup \cite{Biskup2009}. Large fields enter only in this probability estimate, whose proof uses a localized version of the renormalization group. Beyond stability, we also need convergence of the conditioned measures at every fixed number of renormalization group steps. Further, we need estimates for the reconstruction of the fine field from the coarse fields, which control the flowed observables. The localization and fluctuation estimates adapt Combes--Thomas localization \cite{CombesThomas1973}, the Brascamp--Lieb inequality \cite{BrascampLieb1976}, the Bakry--\'Emery criterion \cite{BakryEmery1985}, the Helffer--Sj\"ostrand representation of covariances \cite{HelfferSjostrand1994}, and the Herbst argument as described by Ledoux \cite{Ledoux1999}. The quadratic fluctuation bounds are of Hanson--Wright type \cite{HansonWright1971}, in the form of Adamczak \cite{Adamczak2015}, and the maximal estimates use chaining \cite{Dudley1967,Dirksen2015}.

Compared to the work of Magnen, Rivasseau, and S\'en\'eor \cite{MagnenRivasseauSeneor1993}, our setting is three-dimensional and superrenormalizable. They work in four dimensions, with gauge-fixed Schwinger functions, a fixed infrared cutoff, the trivial topological sector, and counterterms that restore gauge invariance. We work with the Wilson lattice measure without counterterms, and our observables are gauge-invariant flowed Wilson loops. Cao and Chatterjee \cite{CCState2023v4} constructed the state space and proved conditional tightness results; we construct a measure on a version of their state space from the lattice theory. We use their deterministic results on the flow and on the topology of the state space, but not their results for random initial data. The stochastic quantization approach gives local dynamics in three dimensions for the Yang--Mills--Higgs model, including the pure gauge case \cite{ChandraChevyrevHairerShen2024,ChevyrevShen2026Uniqueness}. We construct an equilibrium measure for the pure gauge theory and prove nothing about dynamics. Gross's three-dimensional Abelian result \cite{Gross1983Abelian} and King's construction of the Abelian Higgs model \cite{King1986a,King1986b} are earlier continuum limits of three-dimensional lattice gauge theories; our result concerns a non-Abelian theory.

Our weak-coupling theorems are related to the analysis of Yang--Mills theory on a small torus. The center symmetries fit into the framework of electric and magnetic flux of 't~Hooft \cite{Hooft1979}; we divide out strictly periodic gauge transformations, treat the center operations as symmetries, and include no twist. In the Hamiltonian formulation of the four-dimensional theory on a small spatial torus, L\"uscher \cite{Luscher1983Torus}, L\"uscher and M\"unster \cite{LuscherMunster1984}, and Koller and van Baal \cite{KollerVanBaal1988} showed that the constant modes must be treated nonperturbatively. The leading classical interaction of the constant modes is quartic at the trivial flat connection, and integrating out the other modes adds further terms to the effective Hamiltonian; see van Baal's review \cite{VanBaal2001}. Near regular flat connections our analysis is a version of Laplace's method around a nondegenerate critical manifold, in the sense of Bott \cite{Bott1954CriticalManifolds}. Gonz\'alez-Arroyo, Jurkiewicz, and Korthals Altes \cite{GonzalezArroyo1983}, and Coste and these authors \cite{CosteEtAl1985ZeroMomentum}, studied the zero-momentum modes of the Euclidean lattice theory in a finite box. In this literature the periodic flat connections are often called torons \cite{GonzalezArroyo1998}. Fodor and collaborators combined an exact treatment of the constant modes with perturbation theory for gradient-flow observables on the four-dimensional torus \cite{FodorEtAl2012FiniteVolumeFlow}. The integral of $e^{-Q}$ over the constant modes is the bosonic Yang--Mills integral of three matrices, studied by Krauth, Nicolai, and Staudacher \cite{KrauthNicolaiStaudacher1998}, by Krauth and Staudacher \cite{KrauthStaudacher1998FiniteIntegrals}, and by Austing and Wheater \cite{AustingWheater2001Convergence}. For $\SU$ and three matrices it diverges, and for this reason our theorem at the central triples is stated for compactly supported test functions. Our weak-coupling theorems give rigorous versions of these pictures for the measure we construct, at positive flow time and in the order of limits described above.

\section{Scope}\label{s1:scope}

We record what we do not prove. The limit is taken along the sequences $N_k=nM^k$. We do not compare the limits for different classes of $n$ or different $M$; we expect, but do not prove, that $\lambda$ does not depend on these choices. The torus is fixed: we do not take the infinite-volume limit, and we prove nothing about a mass gap or confinement. The observables are flowed Wilson loops at positive flow time. We do not construct a random distributional connection at time zero, and we do not prove reflection positivity of the limit or the Osterwalder--Schrader axioms. We do not give a rate of convergence in the lattice spacing. The large-field bound and the weak-coupling results are uniform in $\gamma$ below fixed ceilings. The convergence in the main theorem, however, is not shown to be uniform as $\gamma\to0$ or as $\gamma\to\infty$.

\section{Organization}\label{s1:organization}

\Cref{s2:chapter} states our results precisely, and \cref{s3:chapter} gives a detailed outline of the proof. \Cref{ch01:chapter} collects notation and conventions. The proof occupies the remaining sections, grouped into parts. \Cref{part:finite} contains finite-dimensional tools: integration over $\SU$, the block averaging map, reflection positivity, gauge fixing, and the linearized theory. \Cref{part:background} constructs the background fields. \Cref{part:rg} constructs the renormalization group transformation on the small-field event, and proves convergence at each fixed resolution. \Cref{part:reconstruction} estimates the reconstruction of the fine field. \Cref{part:largefield} proves the bound on large fields. \Cref{part:flow} studies the lattice and continuum gradient flows, and \cref{part:stochastic} proves the stochastic estimates for the flowed field. \Cref{part:continuum} constructs the limiting measure and proves the main theorem. \Cref{part:shortdistance,part:weak} prove the results on short distances and on weak coupling.
